% Source printed pp. 324-327 / PDF pp. 332-335.
\section{fppf 拓扑}\label{src:7.3}
\begin{epigraph}
别把时间浪费在空洞的争论上, 就让事情到此为止吧.
\epigraphauthor{William Shakespeare}
\end{epigraph}
尽管本书的两个主要定理并不必需这些内容, 我们仍介绍如下重要事实: 代数空间与代数叠在 fppf 拓扑中分别是层与叠, 且 fppf 群胚与等价关系的商是代数的. 一个结果是, 对任意 (未必光滑的) 代数群, $BG$ 都是代数叠, 例如特征 $p$ 时的 $\mu_p$. fppf 拓扑使我们能在下一节中以适当的一般性定义 gerbe.
\subsection{代数性的 fppf 判据}\label{src:7.3.1-subsection}
\begin{theorem}[代数性的 fppf 判据]\label{src:7.3.1}
\begin{enumerate}[(1)]
\item 若 $X$ 是 $\Sch_{\mathrm{fppf}}$ 上的层, 且存在来自代数空间的 fppf 可表示态射 $U\to X$, 则 $X$ 是代数空间.
\item 若 $\mathcal X$ 是 $\Sch_{\mathrm{fppf}}$ 上的叠, 且存在来自代数空间的 fppf 可表示态射 $U\to\mathcal X$, 则 $\mathcal X$ 是代数叠.
\end{enumerate}
\end{theorem}
\begin{proof}
(1) 由 (2) 得到, 因为等价于层的代数叠是代数空间 (定理~\bookref{5.5.10}). 证明使用类似一般 Keel--Mori 定理的技巧 (见注记~\bookref{5.4.11}). 由于 $U\to\mathcal X$ 可表示, $\mathcal X$ 上的相对 Hilbert 叠 $\mathcal H\!ilb_d$ 是代数叠; 它在 $S\to\mathcal X$ 上的对象是闭子概形 $Z\hookrightarrow U\times_{\mathcal X}S$, 要求在 $S$ 上有限且平坦, 次数为 $d$. 限制到参数化局部完全交 $Z\hookrightarrow U\times_{\mathcal X}S$ 的开子叠 $\mathcal H\!ilb_{d,\mathrm{lci}}$, 可用无穷小提升判据 (\bookref{4.7.1}) 证明 $\coprod_d\mathcal H\!ilb_d\to\mathcal X$ 光滑且满射. 细节见 \cite[定理 6.1]{Art74b}, \cite[定理 10.1]{LMB00} 及 \cite[Tag 06DC]{SP}.
\end{proof}
\begin{remark}\label{src:7.3.2}
若将``fppf''换成``fpqc'', 结论不成立; 见 \cite{MO10}.
\end{remark}
按定义, 代数空间是大平展拓扑中的层, 但它们也是大 fppf 拓扑中的层.
\begin{corollary}[代数叠是 fppf 叠]\label{src:7.3.3}
\begin{enumerate}[(1)]
\item 概形 $S$ 上的代数空间 $X$ 是 $(\Sch/S)_{\mathrm{fppf}}$ 上的层.
\item 概形 $S$ 上的代数叠 $\mathcal X$ 是 $(\Sch/S)_{\mathrm{fppf}}$ 上的叠.
\end{enumerate}
\end{corollary}
\begin{proof}
只需证明 (2). 设 $U\to\mathcal X$ 是光滑呈示, $\mathcal X'$ 是 $\mathcal X$ 在 fppf 拓扑中的叠化. 由于 $U\to\mathcal X'$ 是 fppf 且可表示的, 由定理~\bookref{7.3.1}, $\mathcal X'$ 是代数叠. 由于 $U\times_{\mathcal X}U\cong U\times_{\mathcal X'}U$, 自然映射 $\mathcal X\to\mathcal X'$ 完全忠实. 另一方面, 若 $T\to\mathcal X'$ 是任意来自概形的映射, 纤维积 $U\times_{\mathcal X'}T$ 是 $T$ 上的光滑代数空间. 由于光滑态射在目标上平展局部有截面, 存在平展覆盖 $T'\to T$ 及提升 $T'\to T\to\mathcal X'$ 的映射 $T'\to U$. 复合 $T'\to U\to\mathcal X$ 表明 $\mathcal X'$ 的对象平展局部可提升到 $\mathcal X$, 因而 $\mathcal X\to\mathcal X'$ 本质满. 见 \cite[推论 10.7(a)]{LMB00} 和 \cite[Tag 076V]{SP}.
\end{proof}
\begin{remark}\label{src:7.3.4}
在 Stacks Project \cite{SP} 中, 代数叠按定义是 fppf 拓扑中的叠, 上述结果保证我们的代数叠定义与其定义等价. 代数空间也是 fpqc 拓扑中的层, 这一结果归功于 Gabber; 见 \cite[定理 A.4]{LMB00} 和 \cite[Tag 0APL]{SP}. 还已知具有拟仿射对角的代数叠是 fpqc 拓扑中的叠 \cite[推论 10.7(b)]{LMB00}.
\end{remark}
\begin{exercise}\label{src:7.3.5}
推广命题~\bookref{4.3.6}, 证明代数叠态射的下列性质在目标上 fppf 局部: 可表示, 同构, 开浸入, 闭浸入, 局部闭浸入, 仿射或拟仿射.
\end{exercise}
\subsection{fppf 群胚与商叠}\label{src:7.3.2-subsection}
若 $R\rightrightarrows U$ 是代数空间的 fppf 等价关系, 定义 $U/R$ 为预层 $T\mapsto U(T)/R(T)$ 在大 fppf 拓扑 $\Sch_{\mathrm{fppf}}$ 中的层化. 类似地, 若 $s,t:R\rightrightarrows U$ 是代数空间的 fppf 群胚, 定义 $[U/R]$ 为预叠 $[U/R]^{\mathrm{pre}}$ 在 $\Sch_{\mathrm{fppf}}$ 中的叠化; 其在概形 $T$ 上的纤维范畴 $[U(T)/R(T)]$ 是 $U$ 的 $T$-点范畴, 其中从 $a\in U(T)$ 到 $b\in U(T)$ 的态射是满足 $s(r)=a,t(r)=b$ 的元素 $r\in R(T)$. 由推论~\bookref{7.3.3}, $U/R$ 与 $[U/R]$ 的定义与光滑等价关系或群胚的商 (定义~\bookref{4.4.10}) 相容.
\begin{corollary}\label{src:7.3.6}
\begin{enumerate}[(1)]
\item 若 $R\rightrightarrows U$ 是代数空间的 fppf 等价关系, 则商 $U/R$ 是代数空间.
\item 若 $R\rightrightarrows U$ 是代数空间的 fppf 群胚, 则商 $[U/R]$ 是代数叠.
\end{enumerate}
\end{corollary}
\begin{proof}
由于 $U\to[U/R]$ 是 fppf 且可表示的, 结论由代数性的 fppf 判据 (\bookref{7.3.1}) 得到.
\end{proof}
因此可构造 fppf 群代数空间作用的商叠, 推广定理~\bookref{4.1.10} 中光滑仿射群概形作用的商叠. 首先需要一些定义.
\begin{definition}[主 \texorpdfstring{$G$}{G}-丛]\label{src:7.3.7}
若 $G\to S$ 是 fppf 群代数空间, 则 $S$-概形 $T$ 上的\emph{主 $G$-丛} 是代数空间 $P$, 配备由 $\sigma:G\times_S P\to P$ 给出的 $G$-作用, 使 $P\to X$ 是 $G$-不变的 fppf 态射, 且 $(\sigma,p_2):G\times_S P\to P\times_T P$ 是同构. 主 $G$-丛的态射是 $G$-等变的概形态射. 若存在 $G$-等变同构 $P\cong G\times T$, 则称主 $G$-丛 $P\to T$ 为\emph{平凡的}.
\end{definition}
在上述定义中, 要求 $P$ 是代数空间很重要, 因为我们希望主 $G$-丛满足平展拓扑中的下降. Raynaud 给出了 Abel 簇 $G$ 及不是概形的主 $G$-丛的例子 \cite[XIII 3.2]{Ray70}. 当 $G\to S$ 光滑且仿射时, 这一定义与定义~\bookref{B.1.50} 一致: 事实上, $P$ 是概形, 因为它在 $T$ 上仿射; 每个主 $G$-丛 $P\to T$ 都被光滑覆盖 $P\to T$ 平凡化, 因而可用光滑态射平展局部有截面这一事实, 构造平展覆盖 $T'\to T$, 使 $P_{T'}$ 平凡.
\begin{definition}[商叠与分类叠]\label{src:7.3.8}
设 $G\to S$ 是 fppf 群代数空间, 作用于 $S$ 上的代数空间 $U$. 定义商叠 $[U/G]$ 为 $\Sch/S$ 上的范畴, 其在 $S$-概形 $T$ 上的对象是图
\begin{equation}\label{src:7.3.9}
T\longleftarrow P\longrightarrow U,
\end{equation}
其中 $P\to T$ 是主 $G$-丛, $P\to U$ 是 $G$-等变的概形态射. 态射 $(T'\leftarrow P'\to U)\to(T\leftarrow P\to U)$ 由态射 $T'\to T$ 及 $G$-等变概形态射 $P'\to P$ 组成, 使图
\[
\begin{tikzcd}[ampersand replacement=\&]
P'\arrow[r]\arrow[d]\&P\arrow[r]\arrow[d]\&U\\
T'\arrow[r]\arrow[ur,phantom,"{\square}"]\&T\&
\end{tikzcd}
\]
交换, 且左方块为 Cartesian 方块.
$G$ 的分类叠 $BG$ 定义为商叠 $[S/G]$, 它分类主 $G$-丛 $P\to T$.
\end{definition}
\begin{proposition}\label{src:7.3.10}
若 $G\to S$ 是 fppf 群代数空间, 作用于 $S$ 上的代数空间 $U$, 则商叠 $[U/G]$ 是代数叠. 特别地, 分类叠 $BG$ 是代数的.
\end{proposition}
\begin{proof}
给定对应于对象 \bookeqref{7.3.9} 的映射 $T\to[U/G]$, 有 $\Sch_{\mathrm{fppf}}$ 上叠的 Cartesian 图
\[
\begin{tikzcd}[ampersand replacement=\&]
P\arrow[r]\arrow[d]\&U\arrow[d]\\
T\arrow[r]\arrow[ur,phantom,"{\square}"]\&{[U/G]}
\end{tikzcd}
\]
(见习题~\bookref{3.4.37}). 由于 $P\to T$ 是代数空间的 fppf 态射, $U\to[U/G]$ 是 fppf 可表示态射; 代数性的 fppf 判据 (\bookref{7.3.1}) 表明 $[U/G]$ 是代数叠.
\end{proof}
\begin{exercise}\label{src:7.3.11}
设 $G\to S$ 是 fppf 群代数空间, 作用于 $S$ 上的代数空间 $U$.
\begin{enumerate}[(a)]
\item 证明预叠 $[U/G]^{\mathrm{pre}}$ 在 fppf 拓扑中的叠化是 $[U/G]$. (这推广了习题~\bookref{3.5.21}.)
\item (困难) 给出一个例子, 使 $[U/G]^{\mathrm{pre}}$ 在平展拓扑中的叠化不同构于 $[U/G]$.
\end{enumerate}
\end{exercise}
回顾, $\boldsymbol\mu_{n,\mathbb Z}$ 是 $\mathbb G_{m,\mathbb Z}$ 的子群, 定义为 $\Spec\mathbb Z[x]/(x^n-1)$, 现在可推出 $B\boldsymbol\mu_{n,\mathbb Z}$ 是代数叠. 若 $\Bbbk$ 是特征 $p$ 的域, 则 $\boldsymbol\mu_n:=\boldsymbol\mu_{n,\Bbbk}$ 光滑且 $B\boldsymbol\mu_n$ 是 Deligne--Mumford 叠, 当且仅当 $p$ 不整除 $n$.
\begin{exercise}\label{src:7.3.12}
设 $\Bbbk$ 是域.
\begin{enumerate}[(a)]
\item 给出 $B\boldsymbol\mu_n$ 的显式光滑呈示.
\item 证明 $B\boldsymbol\mu_n$ 等价于 $(\Sch/\Bbbk)_{\text{ét}}$ 上的叠, 其在概形 $T$ 上的对象是对 $(L,\alpha)$, 由 $T$ 上的线丛 $L$ 和平凡化 $\alpha:\mathcal O_T\xrightarrow{\sim}L^{\otimes n}$ 组成.
\item 证明 $B\boldsymbol\mu_n$ 是维数 $0$ 的光滑固有代数叠.
\item 若 $x:\Spec\Bbbk\to B\boldsymbol\mu_n$ 表示典范呈示, 计算切空间 $T_{B\boldsymbol\mu_n,x}$.

\textit{提示:} 答案依赖于特征.
\item (困难) 证明参数化在平展拓扑中可平凡化的 $\mu_n$-丛的叠不是代数的.
\end{enumerate}
\end{exercise}
